\subsection{Normed involutive algebras}

DEFINITION 2. --- A normed involutive algebra is a normed algebra A equipped with an involution \(x \mapsto x^{*}\) such that \(\|x^{*}\| = \|x\|\) for all \(x\) . If A is a Banach algebra, then A is said to be an involutive Banach algebra.

Examples. --- 1) Let X be a locally compact topological space. The Banach algebra \(\mathcal{C}_{b}(\mathrm{X})\) of continuous and bounded complex functions on X, equipped with the norm \(\|f\| = \sup_{x\in\mathrm{X}}|f(x)|\) and with the involution \(f\mapsto\overline{f}\) , is an involutive Banach algebra. The subalgebra \(\mathcal{C}_{0}(\mathrm{X})\) of continuous functions tending to 0 at infinity is a closed involutive subalgebra of \(\mathcal{C}_{b}(\mathrm{X})\) .

\begin{enumerate}
\def\labelenumi{\arabic{enumi})}
\setcounter{enumi}{1}
\item
  Let X be a locally compact topological space and \(\mu\) a positive measure on X. The involutive algebra \(\mathrm{L}^{\infty}(\mathrm{X},\mu)\) (example 2 of I, p.~99) is an involutive Banach algebra, since \(|f| = |\overline{f}|\) for every element \(f \in \mathrm{L}^{\infty}(\mathrm{X},\mu)\) .
\item
  The involutive algebra \(\mathcal{L}(\mathrm{E})\) of continuous endomorphisms of a complex Hilbert space E (I, p.~99, example 3), equipped with the norm
\end{enumerate}

\[
\| u\| = \sup_{\substack{x\in \operatorname{E}\\ \| x\| \leqslant 1}}\| u(x)\|
\]

(EVT, III, p.~14) is an involutive Banach algebra (EVT, V, p.~37, prop. 1).

\begin{enumerate}
\def\labelenumi{\arabic{enumi})}
\setcounter{enumi}{3}
\item
  The involutive algebra \(\mathcal{M}^{1}(G)\) of bounded measures on a locally compact group (I, p.~99, example 4), equipped with the usual norm (example 6 of I, p.~19), is an involutive Banach algebra. Let \(\nu\) be a left Haar measure on G. The involutive Banach algebra \(\mathrm{L}^{1}(\mathrm{G},\nu)\) identifies with a closed subalgebra of \(\mathcal{M}^{1}(\mathrm{G})\) .
\item
  Let \((\mathrm{A}_{i})\) be a family of involutive normed algebras. Let A be the product normed algebra of the \(A_{i}\) (No. 1 of I, p.~15). The algebra A, equipped with the involution \((x_{i})^{*} = (x_{i}^{*})\) , is an involutive normed algebra. If each of the algebras \(A_{i}\) is an involutive Banach algebra, then A is an involutive Banach algebra. We say that A is the product involutive normed algebra (resp. the product involutive Banach algebra) of \(A_{i}\) .
\item
  Let A be a normed involutive algebra, and let \(\widetilde{\mathrm{A}}\) be the normed algebra deduced from A by adjoining a unit element. Equipped with the involution defined in no. 2, the algebra \(\widetilde{\mathrm{A}}\) is a normed involutive algebra. If A is an involutive Banach algebra, then \(\widetilde{\mathrm{A}}\) is also an involutive Banach algebra.
\end{enumerate}

If A is a normed involutive algebra, the closure of an involutive subalgebra is an involutive subalgebra. If M ⊂ A, the smallest closed involutive subalgebra containing M is called the closed involutive subalgebra generated by M; it is the closure of the subalgebra generated by M ∪ M*. If M consists of a single normal element, the closed involutive subalgebra generated by M is commutative, and all its elements are normal.

Similarly, if A is a unital involutive normed algebra and M a subset of A, the smallest closed unital involutive subalgebra containing M is called the closed unital involutive subalgebra generated by M; it is the closure of the unital subalgebra generated by M ∪ M*. If M consists of a single normal element, the closed unital involutive subalgebra generated by M is commutative, and all its elements are normal.

The quotient of an involutive normed algebra by a closed two-sided self-adjoint ideal, the completion, and the opposite of an involutive normed algebra are naturally involutive normed algebras.

If A is a normed involutive algebra, the set \(A_{h}\) of the Hermitian elements of A is a real normed vector space.

Lemma 4. --- Let A be a normed involutive algebra. For every continuous linear form f on A, we have \(\|f^{*}\|=\|f\|\) . If moreover f is Hermitian, then \(\|f\|=\|f|A_{h}\|\) .

The first assertion follows from the definitions. For the second, let \(g\) be the restriction of \(f\) to \(\mathrm{A}_h\) . We have \(\| f \| \geqslant \| g \|\) . Let us prove the reverse inequality. For every \(\varepsilon > 0\) , there exists \(x \in \mathrm{A}\) such that \(\| x \| \leqslant 1\) and \(|f(x)| \geqslant \|f\| - \varepsilon\) . By multiplying x by a complex number of modulus 1, we may assume \(f(x) \geqslant 0\) . Then the element \(\frac{1}{2}(x + x^{*})\) belongs to \(A_{h}\) and has norm \(\leqslant 1\) . We have

\[
\left| g \left(\frac {1}{2} (x + x ^ {*})\right) \right| = \frac {1}{2} | f (x) + f (x ^ {*}) | = f (x) \geqslant \| f \| - \varepsilon
\]

hence \(\|g\|\geqslant\|f\|-\varepsilon\) . From this we deduce that \(\|g\|\geqslant\|f\|\) .

In what follows, we shall identify continuous Hermitian linear forms on A with continuous real linear forms on \(A_{h}\) .

Lemma 5. --- Let A be an involutive Banach algebra.

\begin{enumerate}
\def\labelenumi{\alph{enumi})}
\item
  For every \(x \in A\) , one has \(\exp(x)^{*} = \exp(x^{*})\) ;
\item
  Let \(x \in \mathrm{A}_h\) be a Hermitian element. Then \(\exp(ix)\) is unitary.
\end{enumerate}

Indeed, since the involution on A is continuous, we have

\[
\exp (x) ^ {*} = \left(\sum_ {n = 0} ^ {\infty} \frac {x ^ {n}}{n !}\right) ^ {*} = \sum_ {n = 0} ^ {\infty} \frac {(x ^ {*}) ^ {n}}{n !} = \exp (x ^ {*})
\]

for all \(x \in A\) (formula (9) of I, p.~78). If \(x \in A_{h}\) , it then follows that

\[
\exp (i x) ^ {*} = \exp ((i x) ^ {*}) = \exp (- i x) = \exp (i x) ^ {- 1}
\]

(formula (11) of I, p.~78).

